\section{Good reductions and ordinary base change}\label{bc:section}
\label{sec:base-change}

We retain the finite charts and product presentations of
Proposition~\ref{str:products}.  We will prove ordinary coefficient base
change on suitable reductions, simultaneously for every module on each
chart.  The transverse modules used later need not be flat over the base,
so the proof will keep track of an ordinary complex and its actual chain
homotopies.

Fix a normalized source chart and shrink its normal affine base around a
point of the stratum.  The chart properties in
Proposition~\ref{loc:maximal-charts} give
\[
 h:Z\longrightarrow\Spec R,\qquad
 h_*\cO_Z=\widetilde R,\qquad \dim Z=2n,
\]
where $Z$ is smooth and $h$ is projective and equidimensional of relative
dimension $n$.  Fix a nowhere zero volume form
$\sigma_X\in H^0(X,\omega_X)$; its pullback $\sigma_Z$ is a frame of
$\omega_Z$.  A sufficiently positive embedding over $R$, followed by a
linear projection whose center misses the chosen fiber, gives after
shrinking a finite map
\begin{equation}\label{bc:finite-projection}
 \pi:Z\longrightarrow Y=\PP^n_R.
\end{equation}
Indeed the chosen sections have no common zero on that fiber, hence on
a neighborhood by properness.  The resulting map is quasi-finite because
its defining line bundle is relatively ample, and is therefore finite.
Choose also a finite dominant slice
$\Spec E\to\Spec R$ whose total space is smooth, and put
\[
                 \mathcal V=\pi_*\cO_Z,\qquad
                 E_Y=E\otimes_R\cO_Y.
\]
The finite algebras $\mathcal V$ and $E_Y$ have regular total spaces
$Z$ and $\PP^n_E$.  Quasi-compactness of the stratum lets us make these
choices on a finite collection of source charts.

For a module $E$ over a ring $R$, write $\operatorname{add}_R(E)$ for the
modules isomorphic to direct summands of finite direct sums of $E$.  Our
first target on one chart is a bounded complex $C$ with terms in
$\operatorname{add}_R(E)$ and, for every $R$-module $M$, natural
identifications
\begin{equation}\label{bc:coefficient-target}
 \vartheta_M^i:H^i(C\otimes_R M)
       \xrightarrow{\ \sim\ }H^i(Z,h^*\widetilde M).
\end{equation}
The pullback $h^*$ and all tensor products here are ordinary.
These identifications must respect the specified coefficient maps.  Namely,
for
\[
 \gamma_M^i(C):H^i(C)\otimes_RM\longrightarrow H^i(C\otimes_RM),
 \qquad [z]\otimes m\longmapsto[z\otimes m],
\]
and the usual map $\beta_M^i$ obtained by multiplying a cohomology class
on $Z$ by the pulled-back section $m$, we require
\begin{equation}\label{bc:canonical-compatibility}
 \vartheta_M^i\gamma_M^i(C)
       =\beta_M^i(\vartheta_R^i\otimes\id_M).
\end{equation}
The remaining target in a good reduction is an actual homotopy equivalence,
locally on $\Spec R$, from $C$ to a complex with zero differential.  The
following elementary observation explains why this suffices for all
ordinary modules at once.

\begin{lemma}[The ordinary homotopy criterion]\label{bc:homotopy-criterion}
Let $C$ be a bounded complex of $R$-modules.  If there are chain maps
$C\rightleftarrows D_0$ that are inverse up to chain homotopy and
$d_{D_0}=0$, then $\gamma_M^i(C)$ is an isomorphism for every $R$-module
$M$ and every $i$.  If \eqref{bc:coefficient-target} and
\eqref{bc:canonical-compatibility} hold, the canonical map $\beta_M^i$
is therefore an isomorphism.
\end{lemma}

\begin{proof}
Tensor the two maps and their chain homotopies with $M$.  They remain
homotopy inverses between $C\otimes_RM$ and $D_0\otimes_RM$.  The maps
$\gamma_M^i$ are natural in chain maps, and
$\gamma_M^i(D_0)$ is the identity on $D_0^i\otimes_RM$.  Naturality
under the two homotopy equivalences proves the assertion for $C$.
This uses no flatness of $M$, of $D_0$, or of the terms of $C$.
\end{proof}

We next construct the coefficient complex, then turn a generic homotopy
splitting into the local splitting required by this criterion.

\subsection{Finite smooth covers as coefficient modules}

The following comparison supplies one additive generator for the finite
regular algebras above.  It makes no flatness assertion over their common
normal base.

\begin{proposition}[Comparison of finite regular algebras]\label{bc:finite-modules}
Let $L$ be a complete noetherian normal local domain, and let $A_1,A_2$
be finite dominant local $L$-algebras that are regular local rings.
Put $d_i=\rk_L A_i$.
If these rings have positive characteristic and arise by completing local
rings of varieties over a perfect field, then
\begin{equation}\label{bc:regular-comparison}
                 A_1^{\oplus d_2}\simeq A_2^{\oplus d_1}
                 \qquad\text{as $L$-modules}.
\end{equation}
The same conclusion holds in characteristic zero when the diagram is the
completion of an algebraic finite diagram at a closed point over $\C$.
Finite products of the indicated local algebras can also be compared,
using their total generic ranks.
\end{proposition}

\begin{proof}
Suppose first that the characteristic is $p>0$.  Write $d=\dim L$ and
$\tau=\operatorname{trdeg}_k\kappa(L)$, where $k$ is the perfect ground
field.  All residue fields of the $A_i$ are finite extensions of
$\kappa(L)$, so their $p$-degrees are $p^\tau$.  For $q=p^a$, finite flat
Frobenius on a regular local ring, as in Kunz's theorem
\cite{Kunz1969}, gives
\begin{equation}\label{bc:frobenius-rank}
       F_*^aA_i\simeq A_i^{\oplus q^{d+\tau}}.
\end{equation}
Here the module structure on the left is through $x\mapsto x^q$.
The formula follows either from a regular system of parameters and a
residue-field $q$-basis, or by taking the degree of Frobenius on the
fraction field.  Completion does not change this formula.  In particular,
the exponent is the same for both algebras, including at a nonclosed
point.

The two modules in \eqref{bc:regular-comparison} have the same generic
rank.  Clearing denominators in a generic isomorphism and its inverse
gives maps
\[
 u:A_1^{\oplus d_2}\longrightarrow A_2^{\oplus d_1},\qquad
 v:A_2^{\oplus d_1}\longrightarrow A_1^{\oplus d_2},
 \qquad vu=r\id,\quad uv=r\id
\]
for one nonzero $r\in L$.  We estimate the size of an invertible block of
$F_*^a(r)$ on $F_*^aA_i$, viewed as a free $A_i$-module.  If
$\mathfrak n_i$ is the maximal ideal, its cokernel needs
\[
 \dim_{\kappa(A_i)}
 \bigl(F_*^a(A_i/rA_i)/\mathfrak n_iF_*^a(A_i/rA_i)\bigr)
 =q^\tau\operatorname{length}_{A_i}
       A_i/(r,\mathfrak n_i^{[q]})
\]
generators.  Since $\mathfrak n_i$ has $d$ generators,
$\mathfrak n_i^{d(q-1)+1}\subset\mathfrak n_i^{[q]}$.
The Hilbert--Samuel polynomial of $A_i/rA_i$, of dimension at most $d-1$,
therefore bounds this number by $O(q^{d+\tau-1})$.  The case $d=0$ is
immediate, since $r$ is a unit.  Elementary invertible row and column
operations over $A_i$ now exhibit an identity block of size
$q^{d+\tau}-O(q^{d+\tau-1})$ in $F_*^a(r)$.

Apply $F_*^a$ to $u,v$.  The identity block just found, repeated $d_2$
times, factors through $F_*^a(A_2^{\oplus d_1})$.  Consequently
\[
 A_1^{\oplus d_2(q^{d+\tau}-O(q^{d+\tau-1}))}
 \quad\text{is an $L$-module summand of}\quad
 A_2^{\oplus d_1q^{d+\tau}}.
\]
Finite modules over the complete local ring $L$ have Krull--Schmidt
decompositions.  For each indecomposable isomorphism class $T$, let
$m_i(T)$ be its multiplicity in $A_i$.  The displayed summand relation,
divided by $q^{d+\tau}$ and followed by $a\to\infty$, gives
$d_2m_1(T)\le d_1m_2(T)$.  Interchanging the indices gives equality.
This proves \eqref{bc:regular-comparison}.

Here are details of the characteristic-zero passage.  Henselize the
algebraic base at the chosen point.  The idempotents separating the
finitely many points of a finite algebra descend over an \'{e}tale
neighborhood.  Thus individual factors can be treated in a finite
algebraic diagram; after shrinking that diagram they remain regular and
dominant over a normal integral base.  Let $M=A_1^{\oplus d_2}$ and
$N=A_2^{\oplus d_1}$ denote the corresponding finite modules before
completion.  Choose finite generators $\phi_1,\ldots,\phi_s$ of the
ordinary module $\Hom(M,N)$.  Its completion is
$\Hom(\widehat M,\widehat N)$, because $M$ is finitely presented and
completion is flat.

An isomorphism $\widehat M\simeq\widehat N$ exists precisely when some
linear combination of the $\phi_i$ is surjective modulo the maximal
ideal.  For the forward implication, reduce the coefficients of a
completed isomorphism.  For the converse, Nakayama gives surjectivity,
and a surjective map of the same generic rank between torsion-free
modules has zero kernel.  Existence of the required linear combination
is a finite matrix condition: at least one maximal minor of
$\sum z_i\overline\phi_i$ is a nonzero polynomial in the $z_i$.

Spread the finite presentations, the marked point, and this Hom
calculation over a finitely generated characteristic-zero domain.
Generic freeness applied to the kernels and cokernels of the finite
presentation matrices makes the Hom calculation and its reduction at
the marked point commute with specialization.  Retain normality,
regularity, ranks, and the dimensions of the two residual modules.
If all the maximal-minor polynomials vanished in characteristic zero,
their finitely many coefficients would vanish in every remaining
specialization.  At any good positive-characteristic specialization,
however, the first part supplies an isomorphism of the completed
modules.  This is a contradiction.  A factor that splits further under
specialization causes no difficulty: the first part applies to all its
regular factors and their sum.  Finally, comparisons with one fixed
factor and addition prove the assertion for finite products.
\end{proof}

\begin{remark}\label{bc:residue-degrees}
Lengths in this comparison are always taken over the common target when they
are compared.  For a finite local algebra $A/L$ and a finite-length
$A$-module $Q$,
\[
 \operatorname{length}_L Q
   =[\kappa(A):\kappa(L)]\operatorname{length}_A Q.
\]
Thus tensoring \eqref{bc:regular-comparison} with any fixed finite-length
$L$-module compares the corresponding lengths with their residue-field
degrees included.  Dividing two such comparisons, one for a test module
and one for a fixed reference module of nonzero length, gives the same
ratio on every regular factor.  This is the form needed at generic
points of the support in the Euler calculation.
\end{remark}

Proposition~\ref{bc:finite-modules}, applied to the completed local
factors at closed points, gives
\begin{equation}\label{bc:additive-generator}
              \mathcal V\in\operatorname{add}(E_Y)
              \quad\text{locally on $Y$}.
\end{equation}
To descend the assertion from completion, consider the composition
pairing
\[
 \mathcal Hom(E_Y,\mathcal V)\otimes
 \mathcal Hom(\mathcal V,E_Y)\longrightarrow
 \mathcal End(\mathcal V).
\]
Membership of $\id_{\mathcal V}$ in its image is exactly the existence
of a retraction through a finite sum of $E_Y$'s.  Flat completion
commutes with these finite Hom modules and detects membership by the
cokernel, so the retraction exists locally before completion.  A finite
open cover supplies finitely many retraction witnesses.

Let
\[
 \mathsf H=\End_R(E),\qquad
 \mathsf H_Y=\mathsf H\otimes_R\cO_Y,\qquad
 \mathcal P=\mathcal Hom_Y(E_Y,\mathcal V).
\]
The right action on $\mathcal P$ is by precomposition, and $E$ is a left
$\mathsf H$-module by evaluation.  Flatness of $Y/R$ identifies
$\mathsf H_Y$ with $\mathcal End_Y(E_Y)$.  For a left $\mathsf H$-module
$Q$, set $Q_Y=\cO_Y\otimes_RQ$, with its induced left
$\mathsf H_Y$-action.  Our ordinary sheaf tensor notation means
\[
 \mathcal P\otimes_{\mathsf H}Q
       :=\mathcal P\otimes_{\mathsf H_Y}Q_Y.
\]
The same convention applies to the other right $\mathsf H_Y$-module
sheaves below.  By
\eqref{bc:additive-generator}, $\mathcal P$ is locally a finite
projective right $\mathsf H_Y$-module and evaluation gives
\begin{equation}\label{bc:evaluation}
               \mathcal P\otimes_{\mathsf H_Y}E_Y\simeq\mathcal V.
\end{equation}
These assertions follow first for $E_Y$ and then for sums and summands;
no general Morita equivalence or $R$-flatness of $E$ is being asserted.

\subsection{One twist bound for all coefficient modules}

\begin{lemma}\label{bc:uniform-twists}
Let $R$ be noetherian, let $\mathsf H$ be an algebra finite over its
central subring $R$, and let $\mathcal P$ on $Y=\PP^n_R$ be locally
finite projective as a right $\mathsf H_Y$-module.  There is an integer
$m_0$ such that, for every $m\ge m_0$ and every left $\mathsf H$-module
$Q$,
\begin{align}
 H^j(Y,\mathcal P(m)\otimes_{\mathsf H}Q)&=0
                     &&(j>0),\label{bc:uniform-vanishing}\\
 H^0(Y,\mathcal P(m))\otimes_{\mathsf H}Q
       &\xrightarrow{\ \sim\ }
 H^0(Y,\mathcal P(m)\otimes_{\mathsf H}Q).
                     \label{bc:uniform-sections}
\end{align}
The arrow is the canonical multiplication map, and
$H^0(Y,\mathcal P(m))$ is a finite projective right $\mathsf H$-module.
The bound is independent of $Q$.
\end{lemma}

\begin{proof}
Set $\mathcal K_0=\mathcal P$.  Relative generation on projective space
constructs $n+2$ exact sequences
\begin{equation}\label{bc:finite-resolution}
 0\longrightarrow\mathcal K_{r+1}\longrightarrow\mathcal F_r
   \longrightarrow\mathcal K_r\longrightarrow0,
 \qquad 0\le r\le n+1,
\end{equation}
where each $\mathcal F_r$ is a finite sum of $\mathsf H_Y(-a_{rj})$.
To obtain such a surjection, choose ordinary coherent generators after
twisting and extend their maps $\mathsf H_Y$-linearly.  Since
$\mathcal K_r$ is locally projective over $\mathsf H_Y$, the sequence
locally splits over that algebra.  Induction shows that every kernel
$\mathcal K_{r+1}$ is again locally finite projective over
$\mathsf H_Y$.

The sheaf $\mathsf H_Y$ is flat as a right $\mathsf H$-module: its tensor
functor is $\cO_Y\otimes_R-$, and $Y$ is $R$-flat.  Therefore all the
kernels just constructed are $\mathsf H$-flat, and every sequence in
\eqref{bc:finite-resolution} stays exact after tensoring with an
arbitrary $Q$.  Choose $m_0\ge a_{rj}$ for all the finitely many twists.
For $m\ge m_0$ each $\mathcal F_r(m)\otimes_{\mathsf H}Q$ is a sum of
$\cO_Y(m-a_{rj})\otimes_RQ$.  The coordinate \v{C}ech calculation on
projective space gives zero higher cohomology for these sheaves for
every $R$-module $Q$; their sections are the degree $m-a_{rj}$
homogeneous polynomials with coefficients in $Q$.

For $j>0$, the long exact sequences consequently give
\begin{equation}\label{bc:dimension-shift}
 H^j(Y,\mathcal K_r(m)\otimes_{\mathsf H}Q)
 \simeq
 H^{j+n+2-r}(Y,\mathcal K_{n+2}(m)\otimes_{\mathsf H}Q).
\end{equation}
The $n+1$ standard affine opens of $Y$ have affine intersections, so
quasi-coherent cohomology vanishes in degrees greater than $n$.
Thus the left side is zero for $r=0,1,2$ and $j>0$.
In particular this proves \eqref{bc:uniform-vanishing}.  We have made
no assertion about the projective dimension or the low-degree
cohomology of the final kernel $\mathcal K_{n+2}$.

The first two sequences, together with the vanishings for
$\mathcal K_1$ and $\mathcal K_2$, give
\begin{equation}\label{bc:sections-presentation}
 H^0(Y,\mathcal P(m)\otimes_{\mathsf H}Q)
 =\operatorname{coker}\bigl(
    H^0(Y,\mathcal F_1(m))\otimes_{\mathsf H}Q
       \longrightarrow
    H^0(Y,\mathcal F_0(m))\otimes_{\mathsf H}Q\bigr).
\end{equation}
For $Q=\mathsf H$ this is a finite presentation of
$G_m=H^0(Y,\mathcal P(m))$ by finite free right $\mathsf H$-modules.
Right exactness of tensoring this same presentation proves
\eqref{bc:uniform-sections}, with its canonical map.
Finally, tensoring any short exact sequence of coefficient modules
with $\mathcal P(m)$ remains exact by $\mathsf H$-flatness, and taking
sections remains exact by \eqref{bc:uniform-vanishing}.
Equation~\eqref{bc:uniform-sections} says that $G_m\otimes_{\mathsf H}-$
is exact.  Hence $G_m$ is flat; being finitely presented by
\eqref{bc:sections-presentation}, it is projective.
\end{proof}

Choose an integer $\nu\ge\max(1,m_0)$.  We reserve $\nu$ for the Koszul
exponent; the powers of the stratum ideal keep their notation $I^k$.
The coordinate sections $y_0^\nu,\ldots,y_n^\nu$ have no common zero
on $Y$.  Their Koszul
coresolution is the locally split exact sequence
\begin{equation}\label{bc:koszul}
 0\longrightarrow\mathcal P\longrightarrow
 \mathcal P(\nu)^{\oplus(n+1)}\longrightarrow\cdots\longrightarrow
 \mathcal P((n+1)\nu)\longrightarrow0.
\end{equation}
Local splitting follows by contracting the Koszul complex wherever one
of the coordinate sections is a unit.  Let $K_\nu(\mathcal P)$ be the
complex of sections of its positive-twist terms, beginning in degree
zero.  Lemma~\ref{bc:uniform-twists} makes it a bounded complex of finite
projective right $\mathsf H$-modules representing
$\mathbf R\Gamma(Y,\mathcal P)$.  Set
\[
 C_\nu(\mathcal V)=K_\nu(\mathcal P)\otimes_{\mathsf H}E.
\]
Its terms lie in $\operatorname{add}_R(E)$ and are, explicitly,
$H^0(Y,\mathcal V(j\nu))^{\oplus\binom{n+1}{j}}$ for
$1\le j\le n+1$, in degree $j-1$.
For every $R$-module $M$, ordinary associativity and
\eqref{bc:evaluation} give
\[
 \mathcal P\otimes_{\mathsf H}(E\otimes_RM)
          =\mathcal V\otimes_RM.
\]
Apply Lemma~\ref{bc:uniform-twists} with the left $\mathsf H$-module
$Q=E\otimes_RM$.  The locally split sequence \eqref{bc:koszul} and the
finite projection formula for $\pi$ give natural identifications
\begin{equation}\label{bc:ordinary-complex}
 C_\nu(\mathcal V)\otimes_R M
      \simeq\mathbf R\Gamma(Y,\mathcal V\otimes_RM)
      =\mathbf R\Gamma(Z,h^*\widetilde M).
\end{equation}
The tensor and $h^*$ in this formula are ordinary operations.  No Tor
vanishing for $E$ or $M$ over $R$ has been used.

Here is the canonical-map content of \eqref{bc:ordinary-complex}.  Use
the standard affine cover of $Y$ and its \v{C}ech complex.  The
augmentation from $\mathcal V\otimes_RM$ to the positive Koszul complex
gives a quasi-isomorphism of their \v{C}ech total complexes, since
\eqref{bc:koszul} is locally split before coefficients are tensored.
The maps from global sections of the positive terms to their \v{C}ech
complexes are quasi-isomorphisms by \eqref{bc:uniform-vanishing};
\eqref{bc:uniform-sections} identifies those sections with the terms
of $C_\nu(\mathcal V)\otimes_RM$.  This gives the natural zigzag meant
by \eqref{bc:ordinary-complex}.  The augmentation, \v{C}ech maps,
section identifications, and finite projection formula are all natural
for maps of coefficient modules.
In particular it is natural for $R\to M$, $1\mapsto m$.  On cocycles
that map sends a class to its product with the pulled-back section $m$.
Thus the induced cohomology identifications are precisely
\eqref{bc:coefficient-target} and satisfy
\eqref{bc:canonical-compatibility}.

The same construction applies to any coherent $\mathcal W$ locally in
$\operatorname{add}(E_Y)$, using
$\mathcal Hom_Y(E_Y,\mathcal W)$ in place of $\mathcal P$ and a
sufficiently large exponent.  In the explicit section description,
an $\cO_Y$-linear map of such sheaves induces an actual chain map
$C_\nu(\mathcal W)\to C_\nu(\mathcal W')$ whenever the same $\nu$ is
valid for both.  These maps respect composition.

We also need a stronger comparison between different exponents for one
sheaf.  The bounded projective $\mathsf H$-complexes
$K_\nu(\mathcal P)$ represent the same object
$\mathbf R\Gamma(Y,\mathcal P)$.  Full faithfulness of bounded
projective complexes in the derived category of right $\mathsf H$-modules
therefore lifts its identity to chain maps that are inverse up to
$\mathsf H$-linear chain homotopies.  Tensoring these actual maps and
homotopies with $E$ gives
\begin{equation}\label{bc:power-comparison}
                  C_\nu(\mathcal V)\simeq C_{\nu'}(\mathcal V)
                  \quad(\nu,\nu'\text{ sufficiently large})
\end{equation}
as homotopy equivalences.  The maps and homotopies commute with central
multiplication by elements of $R$.  We used full faithfulness over
$\mathsf H$, before tensoring with $E$, and then retained the resulting
maps and homotopies over $R$.

Fix one such complex $C=C_{\nu_0}(\mathcal V)$ in characteristic zero.
Over $\operatorname{Frac}(R)$ it splits as a complex of vector spaces.
Its cohomology embeds as a summand of a finite zero-differential complex
whose terms are sums of $E\otimes_R\operatorname{Frac}(R)$.  Clearing
denominators in the two maps and in their homotopy gives a nonzero
$g\in R$, a bounded zero-differential complex $D$ with terms finite
sums of $E$, and actual maps and a homotopy
\begin{equation}\label{bc:cleared-homotopy}
 C\xrightarrow{a}D\xrightarrow{b}C,
 \qquad ba-g\id_C=dH_0+H_0d.
\end{equation}
For example, clear separate denominators in the two maps, and then a
further common denominator in the scaled homotopy; rescaling one map
and $g$ gives the displayed identity.  All its terms are finite
module maps.  It is therefore one fixed finite datum that can be
spread out.

The central factor will remain this same $g$ when the exponent changes.
Indeed, let $f:C\to C_{\nu'}(\mathcal V)$ and
$j:C_{\nu'}(\mathcal V)\to C$ be the homotopy inverses induced by the
$\mathsf H$-comparisons, and write
$fj-\id=dJ+Jd$.  Then the actual chain maps $aj$ and $fb$ satisfy
\begin{equation}\label{bc:transported-homotopy}
 (fb)(aj)-g\id_{C_{\nu'}(\mathcal V)}
       =d(fH_0j+gJ)+(fH_0j+gJ)d.
\end{equation}
This follows by inserting \eqref{bc:cleared-homotopy} and using
$fgj=gfj$.  It introduces no additional clearing factor.

\subsection{The differential identity attached to a clearing factor}

Choose a finite Noether normalization
$\mu:\Spec R\to P=\mathbb A^n_\C$, with coordinates $u_1,\ldots,u_n$,
and put
\[
 e_\mu=[\operatorname{Frac}(R):\C(u_1,\ldots,u_n)].
\]
Choose a nonzero $\Delta_\mu\in\C[u_1,\ldots,u_n]$ such that
$R[1/\Delta_\mu]$ is finite \'{e}tale and free of rank $e_\mu$ over
$\C[u_1,\ldots,u_n][1/\Delta_\mu]$, and such that $h$ is smooth over
this locus.  These are fixed finite data.  In particular the retained
rank will later bound the degrees of completed local branches, including
at points outside this open set.

Set $a=\mu h$.
The morphism $a$ is projective, with smooth source of dimension $2n$
and smooth target of dimension $n$.  Write
$du=du_1\wedge\cdots\wedge du_n$.  On the dense locus where $\mu$ is
\'{e}tale and $h$ is smooth, the volume form defines the relative top form
$\eta=\sigma_Z/du$.  Since $h_*\cO_Z=\cO_{\Spec R}$ and
$\omega_Z$ is trivial, the corresponding lowest Hodge sheaf is
$R\eta$, viewed on $P$.

\begin{lemma}\label{bc:differential-identity}
For the fixed $g\ne0$ in \eqref{bc:cleared-homotopy}, on a finite
principal-open cover of $P$ there are finitely many $t_j\in R$ and
algebraic differential operators $D_j$ with
regular coefficients on the respective open sets such that
\begin{equation}\label{bc:differential-equation}
                  \eta=\sum_j D_j(g t_j\eta).
\end{equation}
The identity is in the middle relative de Rham module on the generic
smooth locus, but its operators are regular also at the other points
of the chosen open sets.  For one integer $o_g$ fixed in characteristic
zero, they have expressions
\begin{equation}\label{bc:operator-bound}
 D_j=\sum_{|\alpha|\le o_g}c_{j,\alpha}(u)\partial_u^\alpha,
 \qquad
 \partial_u^\alpha=
 \partial_{u_1}^{\alpha_1}\cdots\partial_{u_n}^{\alpha_n},
\end{equation}
where each $c_{j,\alpha}$ is regular on its specified open set.
\end{lemma}

\begin{proof}
Use the constant polarizable Hodge module on $Z$ in right
$\mathcal D_Z$ notation.  Its underlying module is $\omega_Z$, with
lowest filtration index $-2n$.  The filtered direct-image construction
and projective strictness give
\begin{equation}\label{bc:lowest-filtration}
 F_{-2n}\mathcal H^0(a_+\omega_Z)=a_*\omega_Z.
\end{equation}
To track the index, use the transfer Spencer resolution.  Its term in
degree $-i$ is filtered by pieces with source filtration lowered by
$i$ and by the nonnegative order of target differential operators.
At total index $-2n$, all $i>0$ terms vanish; in degree zero only
order zero remains.  The lowest filtered complex is consequently
$\mathbf Ra_*\omega_Z$.  Strictness identifies its degree-zero
cohomology with the left side of \eqref{bc:lowest-filtration}.
These conventions and strictness are the ones in
\cite[Sections~B.2--B.3 and Theorem~1.4]{Saito2016}, the latter being
\cite[Theorem~5.3.1]{Saito1988}.

The strict-support decomposition is compatible with the filtration.
The sheaf $a_*\omega_Z$ is torsion-free on $P$: multiplication by a
nonzero base function is injective on the dominant integral smooth
source.  Thus its summands on proper supports are zero.  After changing
from right to left modules by $\omega_P^{-1}$, the sheaf $R\eta$ lies
in the full-support part of $\mathcal H^0(a_+\omega_Z)$, which is
generically the $n$-th relative de Rham cohomology.  This part is an
intersection extension of semisimple algebraic local systems, hence a
semisimple regular holonomic $\mathcal D_P$-module with no simple
constituent on a proper support.  Here we use strict support
\cite[Section~1.2]{Saito2016} and algebraic semisimplicity
\cite[Proposition~4.2.5 and Theorem~4.2.6]{Deligne1971}; restricting to an arbitrary small
analytic open would not justify this semisimplicity assertion.

Inside that semisimple module let $U$ and $U_g$ be the
$\mathcal D_P$-submodules generated by $R\eta$ and $gR\eta$.
At the generic point of $P$, the finite domain $R$ becomes the field
$\operatorname{Frac}(R)$, so $g$ is invertible and these restrictions
agree.
The quotient $U/U_g$ is supported properly.  Semisimplicity makes it a
sum of constituents of the full-support module, so it is zero.
Consequently $\eta$ belongs locally to $\mathcal D_P(gR\eta)$.
Writing this membership using finitely many generators gives
\eqref{bc:differential-equation}.  Quasi-compactness gives a finite
principal-open cover and a finite maximum $o_g$ for the orders.
On an open $D(f)$, a local section of $\mu_*\cO_{\Spec R}$ has the form
$r/f^m$ with $r\in R$.  Absorb multiplication by $f^{-m}$ into the
operator acting on $gr\eta$; its coefficients remain regular on $D(f)$
and its order does not increase.  Thus the $t_j$ may be taken in $R$.
Algebraic differential operators on these opens of affine space have
the expressions \eqref{bc:operator-bound}; their finitely many
coefficient denominators are units on the corresponding opens.
The restriction of the identity to the
generic smooth locus involves only algebraic de Rham classes,
connections, and finite-order operations, and thus is finite algebraic
data that can be reduced after localization.
\end{proof}

\subsection{Finite preparations and selection of the reductions}

The preceding constructions give finite data in characteristic zero.
We first retain these data on a model.  After choosing a good reduction,
we will choose its Koszul exponent and its local Cartier units there.
The coefficient modules themselves are not among the data to be spread.

Spread the projective morphism,
the finite charts and their groupoid, the smooth source charts, the
volume form, the ideals of the strata, and the finite \'{e}tale product
presentations.  Include $h,\pi,E,\mu$ on each member of the finite
chart cover, the coherent identity $h_*\cO_Z=R$, the compatible ambient
direct-image identifications of Proposition~\ref{str:direct-images},
the finite Hom and evaluation calculations defining
$\mathsf H,\mathcal P$, the finite local retraction witnesses for
\eqref{bc:additive-generator}, the $n+2$ sequences in
\eqref{bc:finite-resolution} with local splitting witnesses, one exponent
$\nu_0$ and its complex $C_{\nu_0}$, and the maps and homotopy in
\eqref{bc:cleared-homotopy}.  Include also the finite identities
\eqref{bc:differential-equation}, the expressions
\eqref{bc:operator-bound} with their denominators on the specified open
sets, and the finite operator cover itself.  More explicitly, if this
cover is $P=\bigcup_\lambda D(f_\lambda)$, retain an identity
$\sum_\lambda r_\lambda f_\lambda=1$ in $\C[u_1,\ldots,u_n]$.
Its reduction keeps these principal opens covering every fiber.
Include the element $\Delta_\mu$ with the finite free \'{e}tale
rank $e_\mu$.  Localizing a finitely generated integral model makes the kernels,
images, and cokernels used in these presentations flat over the model
and retains their stated base-change maps.  Retain smoothness of the
source charts and slices, normality and geometric integrality of the
chart bases, dimensions, equidimensionality, and the \'{e}tale and finite
assertions on the specified neighborhoods.  Retain also geometric
integrality of the smooth projective total space and smoothness of $h$
over $D(\Delta_\mu)$.  Invert the stabilizer
orders and the finitely many coefficient denominators.  On every retained
fiber, $\Delta_\mu$ remains nonzero, and the algebra on
$D(\Delta_\mu)$ remains finite \'{e}tale and free of rank $e_\mu$.

Retaining the complex $C_{\nu_0}$ includes retaining its coefficient
meaning.  The spread local splitting witnesses make the same $n+2$
sequences \eqref{bc:finite-resolution} locally split over the fiber's
$\mathsf H_Y$, with the same twists $a_{rj}$.  The proof of
Lemma~\ref{bc:uniform-twists} therefore gives the old bound $m_0$ on
every retained fiber, for every coefficient module there.  In
particular $\nu_0\ge m_0$ is still valid.  The first two sequences give
the same finite cokernel presentation \eqref{bc:sections-presentation}
for each $H^0(Y,\mathcal P(j\nu_0))$ after specialization.  Tensoring
that presentation with the residue field identifies the specialized
section module with the section module of the fiber.  Include the
finitely many coordinate multiplication maps as well.  Thus the
specialization of $C_{\nu_0}$ is the fiber complex
$K_{\nu_0}(\mathcal P)\otimes_{\mathsf H}E$, with
\eqref{bc:ordinary-complex} and its canonical-map compatibility.

For the smooth projective total space retain the nowhere zero volume
form and the coherent conditions
\begin{equation}\label{bc:coherent-cup}
 h^2(\cO_X)=h^{2n}(\cO_X)=1,\qquad
 H^2(\cO_X)^{\otimes n}\longrightarrow H^{2n}(\cO_X)
                   \text{ is nonzero}.
\end{equation}
These are finite cohomological conditions, supplied originally by the
IHS Hodge decomposition.  Retain also the relevant coherent and de Rham
base-change identifications.

This preparation remains available after any fixed finite enlargement
of the algebraic and coherent data on a finite-type model, followed by
restriction to any nonempty principal open on which the required
identities and geometric properties remain valid.  Thus this localization
inverts only finitely many elements.  Make those choices before selecting a
characteristic-zero
closed point of the retained model over $\Q$.  Its residue field is a
number field.  Pull the data to its ring of integers with finitely many
elements inverted, so only finitely many further places are excluded by
the model.
The lemma below supplies arbitrarily large remaining Cartier reductions
of this number-field model.  In particular the finite enlargement and
localization occur before that arithmetic choice.  We use only the
retained algebraic and cohomological properties of the chosen
characteristic-zero fiber.

\begin{lemma}\label{bc:cartier-primes}
There are arbitrarily large remaining residue characteristics, at
degree-one places of a finite extension of the number field, for which
\[
                   \operatorname{Car}(F_*\sigma_X)\ne0.
\]
On the geometric special fiber, the map
\begin{equation}\label{bc:cartier-splitting}
 \psi_X:F_*\cO_X\longrightarrow\cO_X,
 \qquad \psi_X(z)=\operatorname{Car}(F_*(z\sigma_X))/\sigma_X
\end{equation}
satisfies $\psi_X(1)=\delta\in k^*$.
\end{lemma}

\begin{proof}
Put $b=b_2(X)$ and choose a prime $\ell>2b$.  After a finite extension,
both the residual representation on a stable lattice of
$H^2_{\mathrm{et}}(X_{\overline K},\Q_\ell)$ modulo $\ell$ and
$\mu_\ell$ are constant.  At a degree-one good place of norm $p$,
geometric Frobenius has an integral trace $a_p$ satisfying
\begin{equation}\label{bc:trace-congruence}
 |a_p|\le bp,\qquad a_p\equiv b\pmod\ell,
                   \qquad p\equiv1\pmod\ell.
\end{equation}
We use integrality and the weight-two Weil bound, and the equality of
the \'{e}tale and crystalline characteristic polynomials
\cite[Theorem~1]{KatzMessing1974}, using its formulation and the
crystalline applicability stated in
\cite[Chapter~VI, Theorem~3.9 and Section~4]{MilneTannakian2025}.

Discard the finitely many ramified places of the number field and
the places where relative de Rham cohomology or its base change has
torsion.  At a degree-one remaining place the local base is $\Z_p$.
Crystalline--de Rham comparison for the smooth proper lift identifies
the reduction of the crystalline trace with the trace of Frobenius on
$H^2_{\mathrm{dR}}(X_{\F_p})$
\cite[Corollaries~3.8 and~3.10]{BhattDeJong2011}.
Modulo $p$, Frobenius on the de Rham complex factors as its projection
onto $\cO_X$ followed by the map taking functions to their $p$-th
powers in degree zero.  On cohomology write this factorization as
$\beta\alpha$ through $H^2(\cO_X)$.  The reverse composite
$\alpha\beta$ is coherent Frobenius.  The identity
$\tr(\beta\alpha)=\tr(\alpha\beta)$ shows that zero coherent
Frobenius forces $p\mid a_p$.

In that case $a_p/p$ is an integer in $[-b,b]$ congruent to $b$ modulo
$\ell$.  Since $\ell>2b$, it equals $b$, so $a_p=bp$.
If coherent Frobenius vanished at all but finitely many degree-one
good places, the semisimplification of
$H^2_{\mathrm{et}}(X_{\overline K},\Q_\ell)$ would therefore have the
same character as $\Q_\ell(-1)^{\oplus b}$.  Indeed degree-one places
have density one among primes of a number field; the others have norm
at least the square of the underlying rational prime and density zero.
Chebotarev for number fields and the characteristic-zero character
criterion \cite[Theorem~3.2 and Section~5.1.1.3]{Serre2014} give the
assertion: finite-quotient Frobenius density and continuity identify
the two trace functions in the projective limit.

This contradicts Hodge--Tate comparison at a place above $\ell$:
$h^2(\cO_X)=1$ contributes a weight-zero summand, whereas all Tate
summands just found have weight one, in the convention where
$\Q_\ell(-1)$ has weight one.  Comparison applies to the smooth
proper variety over that local field
\cite[Theorem~1.6 and Corollary~1.8]{Scholze2013}; de Rham representations are stable
under subquotients and their filtered dimensions are additive in
exact sequences \cite[Theorem~5.2.1(2) and Propositions~6.3.2--6.3.3]{BrinonConrad2009},
so passing to semisimplification has not lost this obstruction.  Thus coherent Frobenius is nonzero at infinitely many,
and hence arbitrarily large, good residue characteristics.

By \eqref{bc:coherent-cup}, its $n$-fold product is nonzero on top
coherent cohomology.  Serre duality identifies the transpose of this
map with Cartier trace on global top forms.  The space of global top
forms is the frame line and $X$ is geometrically integral and proper,
so \eqref{bc:cartier-splitting} sends $1$ to a nonzero constant.
\end{proof}

All subsequent positive-characteristic assertions are made in these
geometric special fibers.  Choose $p$ larger than every retained
Noether-normalization degree $e_\mu$ and every differential order bound
$o_g$ on the finite chart cover.  \'{E}tale compatibility of Cartier
trace gives the same constant $\delta$ on the source charts.  Thus
$\psi_Z:F_*\cO_Z\to\cO_Z$ pushes down, using $h_*\cO_Z=R$, to
\[
                  \psi_h:F_*R\longrightarrow R,
                  \qquad\psi_h(1)=\delta.
\]

\subsection{Cartier detects every cleared factor}

\begin{lemma}\label{bc:clearing-test}
For every closed point $b\in\Spec R$ in a retained chart, there is a
function $t$ regular near $b$ such that
\begin{equation}\label{bc:unit-test}
                         \psi_h(gt)(b)\ne0.
\end{equation}
\end{lemma}

\begin{proof}
Fix $b$ and one of the principal opens in
Lemma~\ref{bc:differential-identity} containing $\mu(b)$.  In this proof
$t_j,D_j$ denote its data, so the $t_j$ are regular at $b$ and the
coefficients of the $D_j$ are regular at $\mu(b)$.

First work over $K_0=\operatorname{Frac}(R)$, where the family
$Z_0=Z\times_RK_0$ is smooth and proper.  The retained \'{e}tale locus
of $\mu$ identifies $du$ with a basis of $\omega_{K_0/k}$ and gives
the relative frame $\eta=\sigma_Z/du$.  Serre duality gives dual lines
\[
 V_0=H^n(Z_0,\cO_{Z_0})=K_0v,\qquad
 W_0=H^0(Z_0,\omega_{Z_0/K_0})=K_0\eta,\qquad
 \langle v,\eta\rangle=1.
\]
Put
\[
 Z_0^{(p)}=Z_0\times_{\Spec K_0,F_{K_0}}\Spec K_0,\qquad
 V_0^{(p)}=V_0\otimes_{K_0,F_{K_0}}K_0,\qquad
 W_0^{(p)}=W_0\otimes_{K_0,F_{K_0}}K_0.
\]
Flat field base change identifies these twisted vector spaces with
$H^n(Z_0^{(p)},\cO_{Z_0^{(p)}})$ and
$H^0(Z_0^{(p)},\omega_{Z_0^{(p)}/K_0})$, respectively.
Write $v^{(p)}=v\otimes1$ and $\eta^{(p)}=\eta\otimes1$.
They remain dual under the base-changed Serre pairing.
Relative Frobenius $F_{\mathrm{rel}}:Z_0\to Z_0^{(p)}$ and its
finite dual trace have the following $K_0$-linear maps on these lines:
\begin{equation}\label{bc:relative-trace-lines}
 \begin{aligned}
 F_{\mathrm{rel}}^*:V_0^{(p)}&\longrightarrow V_0,
       &v^{(p)}&\longmapsto s\,v,\\
 \operatorname{Tr}_{F_{\mathrm{rel}}}:W_0&\longrightarrow W_0^{(p)},
       &\eta&\longmapsto s\,\eta^{(p)}.
 \end{aligned}
\end{equation}
Here the second map is induced by
$F_{\mathrm{rel}*}\omega_{Z_0/K_0}\to\omega_{Z_0^{(p)}/K_0}$.
It is the Serre-dual transpose of the first, which proves that their
scalars are the same $s\in K_0$.

Taking $p$-th powers of a \v{C}ech cocycle for $v$ gives a class
$\lambda$ in middle relative de Rham cohomology.  Its projection to
$H^n(Z_0,\cO_{Z_0})$ is $F_{\mathrm{rel}}^*v^{(p)}$.  Pairing
$\lambda$ with the top form $\eta$, with the order that gives Serre
evaluation, therefore gives
$\langle\lambda,\eta\rangle=s$.  This distinguishes the coherent
Frobenius map in \eqref{bc:relative-trace-lines} from its de Rham lift.
The lift is horizontal: its $p$-power representatives are closed even
for absolute differentials over $k$, so the connecting operation defining
the Gauss--Manin derivative annihilates them.  The de Rham pairing is
compatible with the connection.  Pairing the reduced identity
\eqref{bc:differential-equation} with this horizontal class gives
\begin{equation}\label{bc:scalar-differential}
                           s=\sum_jD_j(gt_js).
\end{equation}
The pairing, connection, and \eqref{bc:differential-equation} on the
smooth generic family are reductions of their finite characteristic-zero
descriptions.  The scalar identity \eqref{bc:scalar-differential} is
obtained in this fiber from the newly constructed horizontal class.

Factor absolute Frobenius as
$Z_0\xrightarrow{F_{\mathrm{rel}}}Z_0^{(p)}
 \xrightarrow{w}Z_0$, where $w$ is the base change of $F_{K_0}$.
This is the relative Frobenius factorization of
\cite[Tag 0CC6]{StacksProject}.
The relative trace in \eqref{bc:relative-trace-lines} sends
$m\eta$ to $sm\eta^{(p)}$.  The trace for $w$, by flat base change and
the projection formula on top forms, applies the Cartier map
$\operatorname{Car}_{K_0}:F_*\omega_{K_0/k}\to\omega_{K_0/k}$
to its remaining base form.  Transitivity of finite trace therefore
gives the rational identity below; compatibility of trace with this
flat base change is \cite[Tag 0E5L]{StacksProject}.
\begin{equation}\label{bc:relative-cartier}
 \psi_h(m)\,du
       =\operatorname{Car}_{K_0}\bigl(F_*(s m\,du)\bigr)
       \qquad(m\in K_0).
\end{equation}
We suppress $F_*$ in the input notation for $\psi_h$.  The scalar
$s$ stays inside the base Cartier operator.  At this point it is only
a rational element of $K_0$.

Suppose \eqref{bc:unit-test} fails, so
$\psi_h(gR_b)\subset\mathfrak m_{R_b}$.  Complete the regular base at
$\mu(b)$:
\[
 A=k[[v_1,\ldots,v_n]],\qquad v_i=u_i-u_i(\mu(b)),\qquad
 K=\operatorname{Frac}(A),\qquad \mathfrak m=(v_1,\ldots,v_n).
\]
The finite completed algebra $R\otimes_{k[u]}A$ is a product of the
completed local rings at the points above $\mu(b)$.  They are normal
complete local domains; keep the factor $L$ at $b$ and put
$K'=\operatorname{Frac}(L)$.  After tensoring with $K$, the whole
product is finite \'{e}tale of total dimension $e_\mu$: the nonzero
$\Delta_\mu$ is invertible in $K$, and its free \'{e}tale rank was
retained.  Thus its factors are finite separable fields and
\begin{equation}\label{bc:completed-degree}
       [K':K]\le e_\mu<p.
\end{equation}
Write $T=\operatorname{Tr}_{K'/K}$.

We next transport \eqref{bc:relative-cartier} to these fields; this step
does not require $s$ to be regular in $L$.  Finite Frobenius and the
cofinality of the ordinary and Frobenius adic filtrations give
$F_*R_b\otimes_{R_b}L\simeq F_*L$, by
$(F_*r)\otimes x\mapsto F_*(rx^p)$.  Hence $\psi_h$ extends to $L$
and then to a $K'$-linear map $F_*K'\to K'$.  The monomials
$v^\alpha$, $0\le\alpha_i<p$, are a $p$-basis of $K$: decompose power
series according to their exponents modulo $p$ and use perfection of
$k$.  They remain a $p$-basis of the finite separable extension $K'$.
They are also a $p$-basis of the embedded $K_0$, since $K_0$ is finite
separable over $k(u_1,\ldots,u_n)$ and translation by constants of
$k$ changes $u_i$ to $v_i$.  The Cartier maps computed in these bases
agree on $K_0$.  The extended $\psi_h$ and
$m\mapsto\operatorname{Car}_{K'}(F_*(sm\,du))/du$ agree on the
$p$-basis monomials by \eqref{bc:relative-cartier}; they are both
$K'$-linear maps on $F_*K'$, so they agree everywhere.  The same
uniqueness of separable extension shows that the extended derivations
in \eqref{bc:scalar-differential} agree on $K_0$, so that identity also
holds in $K'$ with these operators.

The failed containment extends to
$\psi_h(gL)\subset\mathfrak m_L$.  For every base monomial $a=v^\beta$
and each $t_j$ it follows that
\begin{equation}\label{bc:trace-in-maximal}
                       T(\psi_h(gat_j))\in\mathfrak m.
\end{equation}
Indeed $T(L)\subset A$ by integrality and normality.  An element of
$\mathfrak m_L$ has all its conjugates in the maximal ideal above
$\mathfrak m$: the integral closure in a finite normal closure is
finite over the complete regular ring $A$ and is local by henselianity.
Their sum consequently
lies in $\mathfrak m$.

The common $p$-basis also proves that Cartier commutes with this rational
trace.  If $x=\sum v^\alpha c_\alpha^p$ in $K'$, then
$T(x)=\sum v^\alpha T(c_\alpha)^p$, since $T(c^p)=T(c)^p$.
Cartier for $du=dv_1\wedge\cdots\wedge dv_n$ extracts the coefficient
with $\alpha=(p-1)\mathbf1$ in either field.  Hence
\[
 T\bigl(\operatorname{Car}_{K'}(F_*(x\,du))/du\bigr)
       =\operatorname{Car}_{K}(F_*(T(x)\,du))/du.
\]
Derivations likewise commute with $T$, by uniqueness of their separable
extensions.  Combining this identity with
\eqref{bc:relative-cartier} on $K'$ and
\eqref{bc:trace-in-maximal}, and setting $z_j=T(gt_js)\in K$, gives
\[
 \operatorname{Car}_{K}(F_*(a z_j\,du))/du\in\mathfrak m
                         \quad\text{for every base monomial }a.
\]
Write uniquely
\[
 z_j=\sum_{0\le\alpha_i<p}v^\alpha c_{j,\alpha}^{p},
                         \qquad c_{j,\alpha}\in K.
\]
Multiplication by $v^{(p-1)\mathbf1-\alpha}$ followed by Cartier
extracts exactly $c_{j,\alpha}$.  For another exponent $\beta$ the
exponent vector lies componentwise between $0$ and $2p-2$ and has
Cartier residue $(p-1)\mathbf1$ only when $\beta=\alpha$.
Consequently every $c_{j,\alpha}$ belongs to $\mathfrak m$.  This
proves the required regularity of the traced products, namely
\begin{equation}\label{bc:frobenius-ideal}
                    z_j\in\mathfrak m^{[p]}
                           =(v_1^p,\ldots,v_n^p)A.
\end{equation}
No regular extension of $s$ or $\eta$ across $b$ has entered this
argument.

The reduced expressions \eqref{bc:operator-bound} have coefficients in
$A$ at $\mu(b)$ and ordinary differential order at most $o_g<p$.
They preserve $\mathfrak m^{[p]}$, since the ordinary partial
derivatives annihilate every $v_i^p$.  Their coefficients lie on the
regular base $P$, so they pass through the $K$-linear trace $T$.
Taking $T$ in \eqref{bc:scalar-differential} and using
\eqref{bc:frobenius-ideal} gives $T(s)\in\mathfrak m^{[p]}$.
Cartier consequently sends $T(s)du$ into $\mathfrak m\,du$.
But \eqref{bc:relative-cartier} with $m=1$ gives
\[
 \operatorname{Car}_{K}(F_*(T(s)du))/du
       =T(\psi_h(1))=[K':K]\delta\ne0\pmod{\mathfrak m},
\]
where \eqref{bc:completed-degree} and $\delta\in k^*$ give the final
inequality.  This contradiction proves the lemma.
\end{proof}

\subsection{Frobenius on the actual Koszul complexes}

In a retained reduction, $F_{Y*}\mathcal V$ again belongs locally to
$\operatorname{add}(E_Y)$.  In fact
$F_{Y*}\pi_*\cO_Z=\pi_*F_{Z*}\cO_Z$, and
$F_{Z*}\cO_Z$ is locally free over the smooth source.  Over the
semilocal finite algebra above a target point it is finite projective; hence its
pushforward is locally in $\operatorname{add}(\mathcal V)$, and
\eqref{bc:additive-generator} applies.  Also $F_{R*}E$ is a finite
projective $E$-module when $E$ acts through its Frobenius, so
\begin{equation}\label{bc:frobenius-add}
             F_{R*}\operatorname{add}_R(E)
                             \subset\operatorname{add}_R(E).
\end{equation}
The old bound $m_0$ remains valid for $\mathcal V$ by the spread split
sequences.  Apply Lemma~\ref{bc:uniform-twists} also to
$F_{Y*}\mathcal V$ in this characteristic, and choose one exponent
$\nu$ valid for both sheaves.  This new $\nu$ may depend on the reduction
and on the chart.

There is a literal identification of complexes
\begin{equation}\label{bc:frobenius-koszul}
 F_{R*}C_{p\nu}(\mathcal V)=C_\nu(F_{Y*}\mathcal V).
\end{equation}
Indeed the term indexed by $j$ is identified by finite Frobenius
projection formula as
\[
 F_{R*}H^0(Y,\mathcal V(jp\nu))
      =H^0(Y,(F_{Y*}\mathcal V)(j\nu)).
\]
Multiplication by $y_i^{p\nu}$ before pushforward is multiplication by
$y_i^\nu$ in the pushed module structure.  The coordinate ordering and
Koszul signs agree, so \eqref{bc:frobenius-koszul} identifies the
differentials as well as the terms.

Write $\mathsf m_g$ for multiplication by $g$ on the unpushed sheaf or
complex.  Its pushforward is the map
$F_*\mathsf m_g:F_*z\mapsto F_*(gz)$; the scalar action of $g$ on the
pushed module instead sends $F_*z$ to $F_*(g^pz)$.  The identification
\eqref{bc:frobenius-koszul} respects the former map.

The retained $\nu_0$ and $p\nu$ both satisfy the fiber's twist bound.
Use the $\mathsf H$-projective comparison before tensoring with $E$ to
apply \eqref{bc:transported-homotopy} with $\nu'=p\nu$.  Denote its
maps by $A'=aj:C_{p\nu}(\mathcal V)\to D$ and
$B'=fb:D\to C_{p\nu}(\mathcal V)$, and its homotopy by $H'$.
Applying the additive functor $F_{R*}$ to the actual identity gives
\begin{equation}\label{bc:pushed-homotopy}
 (F_{R*}B')(F_{R*}A')-F_{R*}\mathsf m_g
             =d(F_{R*}H')+(F_{R*}H')d.
\end{equation}
Under \eqref{bc:frobenius-koszul}, this factors the pushed map
$F_*\mathsf m_g$ on $C_\nu(F_{Y*}\mathcal V)$ up to homotopy through
the zero-differential complex $F_{R*}D$.  The central factor in
\eqref{bc:transported-homotopy} was the original $g$.

Fix a closed point $b$, and choose $t$ as in
Lemma~\ref{bc:clearing-test}.  On a neighborhood of $b$, put
$u=\psi_h(gt)\in R^*$.  There are $\cO_Y$-linear maps
\begin{equation}\label{bc:cartier-composite}
 \mathcal V\xrightarrow{\iota}F_{Y*}\mathcal V
 \xrightarrow{F_*\mathsf m_g}F_{Y*}\mathcal V
 \xrightarrow{\rho_t}\mathcal V.
\end{equation}
Here $\iota(z)=F_*(z^p)$ and
$\rho_t(F_*z)=\psi_Z(h^*t\,z)$.
Their composite is $u\id_{\mathcal V}$, because
$\psi_Z(h^*(tg)z^p)=z\,h^*\psi_h(tg)$.
Use the same exponent $\nu$ in the two Koszul coresolutions.  Naturality
makes \eqref{bc:cartier-composite} actual chain maps of their section
complexes, with exactly the same composite.  Compose those chain maps
with \eqref{bc:pushed-homotopy} and divide the last map by $u$.  This
gives actual maps
\[
 C_\nu(\mathcal V)\xrightarrow{\alpha}F_{R*}D
       \xrightarrow{\beta}C_\nu(\mathcal V),
 \qquad \beta\alpha\simeq\id_{C_\nu(\mathcal V)}.
\]
The displayed homotopy is obtained by composing the actual homotopy in
\eqref{bc:pushed-homotopy} with the two chain maps.

For completeness, such a homotopy retract is homotopy split.  If
$C\xrightarrow{\alpha}D\xrightarrow{\beta}C$ has
$\beta\alpha\simeq\id_C$ and $d_D=0$, then
$e=\alpha\beta$ satisfies $e^2\simeq e$.  A homotopy between maps
from $D$ to itself is zero, so $e^2=e$ as an actual degreewise map.
Its images form a zero-differential complex $D'$.  Define
\[
 \alpha'=e\alpha:C\longrightarrow D',\qquad
 \beta'=\beta|_{D'}:D'\longrightarrow C.
\]
Then $\alpha'\beta'=\id_{D'}$ and
$\beta'\alpha'=(\beta\alpha)^2\simeq\id_C$, so these maps are
homotopy inverses.  The projection in $\alpha'$ is necessary:
$\alpha$ itself need not factor through $D'$.  Idempotents split in
modules, so no additional categorical completeness is required.

\begin{proposition}[Ordinary coefficient base change]\label{bc:ordinary}
Perform the finite characteristic-zero preparations above on a
finite-type integral model.  After any fixed finite enlargement of its
algebraic and coherent data, and restriction to any nonempty principal
open retaining the required identities and geometric properties, one can
choose a characteristic-zero closed point and a number-field model with
geometric reductions in arbitrarily large residue characteristics $p$
that retain these data and satisfy
\[
 \psi_X(1)=\delta\in k^*,\qquad
 p>e_\mu,\quad p>o_g
       \quad\text{on every member of the finite chart cover}.
\]
On each such reduction the following assertions hold simultaneously.
On every retained source chart, for every
$R$-module $M$ and every $i$, the canonical map
\begin{equation}\label{bc:canonical-base-change}
 (R^ih_*\cO_Z)\otimes_RM\longrightarrow
                          R^ih_*(h^*\widetilde M)
\end{equation}
is an isomorphism, with ordinary tensor and pullback.  In particular,
for every coherent module $M$ supported on any finite thickening of
$\cS$ in $\widehat{\mathcal B}$,
\begin{equation}\label{bc:ordinary-euler}
 \chi(X,f^*M)=\sum_{i=0}^n(-1)^i
                \chi(\widehat{\mathcal B},N_i\otimes M).
\end{equation}
One reduction works for all these $M$ and all thickening orders.
\end{proposition}

\begin{proof}
The preceding argument makes $C_\nu(\mathcal V)$ homotopy equivalent
to a zero-differential complex near every closed point.  Its
coefficient identifications satisfy
\eqref{bc:coefficient-target} and
\eqref{bc:canonical-compatibility}.  Lemma~\ref{bc:homotopy-criterion}
therefore gives the canonical map \eqref{bc:canonical-base-change} for
every module on each such neighborhood.  Since the base is affine, the
cohomology modules correspond to the quasi-coherent direct-image sheaves;
the map is multiplication by the pulled-back coefficient.

The local homotopy splitting extends to a neighborhood of each closed
point: its maps involve finite modules, the local function $t$, and
the inverse of the unit $u$.  These neighborhoods are chosen before
any coefficient module $M$.  They cover the finite-type chart, since any nonempty
closed complement would have a closed point.  Consequently the local
statement holds as a sheaf statement for every module.  The maps are
compatible with the \'{e}tale overlaps and deck actions because they are
canonical multiplication maps.  Proposition~\ref{str:products}
therefore descends them to every supported thickening.  The ordinary
supported pullback on $X$ is exactly the one constructed there.
Taking the Euler characteristic of Leray and using the spread
chartwise direct-image isomorphisms of
Proposition~\ref{str:direct-images} gives \eqref{bc:ordinary-euler}.
Here $N_i\otimes M$ has the ordinary bounded-support meaning fixed in
Section~\ref{sec:stratum}; all supports are proper.

The only choices required before reduction were the finite diagrams,
resolutions, one cleared homotopy per chart, and bounded differential
identities.  The proof of Lemma~\ref{bc:uniform-twists} works for all
coefficient modules in each retained fiber, with the old bound for
$C_{\nu_0}$ retained by the split sequences.  The exponent $\nu$, the
comparison at $p\nu$, and the local functions $t$ were constructed anew
in that fiber.  No spreading of infinitely many Frobenius identities
or of the individual modules $M$ occurs.  A finite chart cover and
Lemma~\ref{bc:cartier-primes} give the asserted simultaneous choice of
reductions.
\end{proof}
